\documentclass{beamer}
\usetheme{desy}
\usepackage{amsfonts,amsmath,oldgerm,algorithm,algpseudocode}
\usepackage[font=small,labelfont=bf]{caption} % Required for specifying captions to tables and figures

\newcommand{\hrefcol}[2]{\textcolor{uihteal}{\href{#1}{#2}}}
\newcommand{\testcolor}[1]{\colorbox{#1}{\textcolor{#1}{test}}~\texttt{#1}}

% Please see Section 18.1 of Beamer User Guide for all the options \usefonttheme provides
\usefonttheme[onlymath]{serif}
% \usefonttheme{serif} % use this if you would like Serif font throughout (and not just for math)


\title{Get started with the DESY-SINTEF Beamer Theme}
\subtitle{Using \LaTeX\ to prepare slides}
% This can be adjusted accordingly for longer titles
\setlength{\titleboxwidth}{0.55\textwidth}

\author{\href{mailto:saverio.monaco@desy.de}{Saverio Monaco}}
\date{\today}

\begin{document}
\maketitle

\begin{frame}[fragile]{Beamer for UIC presentations}
	\framesubtitle{A quick start}
	If you would like \LaTeX\ in your presentation, Beamer is a great way to go!
	\begin{itemize}
		\item Beamer has a detailed
		      \hrefcol{https://www.ctan.org/tex-archive/macros/latex/contrib/beamer/doc/beameruserguide.pdf}{user
			      manual}, but we will go over the most common features.
		\item This template is designed for a 16:9 aspect ratio, which is the default in PowerPoint and the most common amongst projectors.
		\item The most common of all slide types involve bulleted points, like these. \pause
		      \begin{itemize}
			      \item Placing \verb|\pause| after items will allow you to sequentially unroll points.
			      \item Changing \verb|\begin{itemize}| to \verb|\begin{itemize}[<+->]| achieves a similar effect, although \verb|\pause| gives finer control on the unrolling.
		      \end{itemize}
	\end{itemize}
\end{frame}

\begin{chapter}[plots/min_fuel.png]{0.8}{3}{0}{Slide Layouts}
 \textit{Some standard slide layouts}
\end{chapter}

\begin{frame}[fragile]{Images, columns and tips for good slides}
	\framesubtitle{\textbf{Do use} this optional subtitle where appropriate}
	\begin{columns}
		\begin{column}{0.7\textwidth}
			\begin{itemize}
				\item \textbf{Do keep} a good balance between text and figures.
				      \begin{itemize}
					      \item \textbf{Do use} nested lists wherever appropriate.
				      \end{itemize}
				\item \textbf{Do use} the \verb|\underbrace| and \verb|\overbrace| commands to help explain your equations:
				      \begin{equation*}
					      \begin{aligned}
						      \max_{\lambda, \nu}. \quad & \underbrace{\inf_x \overbrace{L(x, \lambda, \nu)}^\text{Lagrangian}}_\text{Lagrange Dual Function, $g(\lambda, \nu)$} \\
						      \textrm{s.t.} \quad        & \lambda \geq 0
					      \end{aligned}
					      % \max_{\lambda, \nu}. \underbrace{\inf_x \overbrace{L(x, \lambda, \nu)}^\text{Lagrangian}}_\text{Lagrange Dual Function, $g(\lambda, \nu)$}, \quad \textrm{s.t.} \quad \lambda \geq 0
				      \end{equation*}
			\end{itemize}


		\end{column}
		\begin{column}{0.3\textwidth}
			\includegraphics[width=\textwidth]{plots/min_fuel.png}
			\captionof{figure}{Minimum fuel trajectory}
		\end{column}
	\end{columns}
\end{frame}

\themecolor{dark}
\begin{frame}
	\frametitle{A standout frame with title and subtitle}
	\framesubtitle{Subtitle here}
	\centering\large
	A \textbf{\itshape\scshape standout} frame can be used to focus attention
\end{frame}
\themecolor{light}

\begin{frame}[fragile]{More with bullet points}
	\framesubtitle{Alerts and repeats}
	Sections 12.1 through 12.3 of the \hrefcol{https://www.ctan.org/tex-archive/macros/latex/contrib/beamer/doc/beameruserguide.pdf}{Beamer user
		manual} demonstrate many more features, like alerts and repeats.
	\begin{itemize}
		\item The \verb|\alert{}| feature can be particularly useful.
		\item Like \alert<2>{this}.
		\item The highlighting may span multiple slides in a frame.
		      \begin{itemize}
			      \item Previous one did not, but \alert<3->{this one} will.
		      \end{itemize}
		\item This is accomplished by the \alert<4>{hyphen} (i.e. \texttt{-}) in the \verb|\alert<3->| command.
	\end{itemize}
\end{frame}

\begin{frame}{Another images/columns example}
	\begin{columns}
		\begin{column}{0.35\textwidth}
			\includegraphics[width=\textwidth]{plots/min_fuel.png}
			\captionof{figure}{Minimum fuel trajectory}
		\end{column}
		\begin{column}{0.3\textwidth}
			\textbf{Left:} Solving for optimal fuel consumption \textbf{Right:} Solving for optimal time taken
		\end{column}
		\begin{column}{0.35\textwidth}
			\includegraphics[width=\textwidth]{plots/min_time.png}
			\captionof{figure}{Minimum time trajectory}
		\end{column}
	\end{columns}
\end{frame}


\begin{frame}{Blocks}
	Some content will just look better in blocks.
	\begin{columns}
		\begin{column}{0.4\textwidth}
			\begin{block}{Block title}
				This is a \texttt{block}. Its color will match the color of the footline. The \texttt{block} environment is native to Beamer.
			\end{block}
		\end{column}
		\begin{column}{0.6\textwidth}
			\begin{colorblock}[black]{chicagoblue}{Block title}
				This Beamer theme also provides a \texttt{colorblock} environment which is a wrapper on top of the usual \texttt{block} environment in Beamer. It allows you to specify a custom background and font color.
			\end{colorblock}
		\end{column}
	\end{columns}

	\vspace{3ex}
	Here's an example usage:

	\begin{columns}
		\begin{column}{0.5\textwidth}
			\begin{block}{The discrete case}
				...
			\end{block}
		\end{column}
		\begin{column}{0.5\textwidth}
			\begin{colorblock}[black]{chicagoblue}{The continuous case}
				...
			\end{colorblock}
		\end{column}
	\end{columns}
\end{frame}


\begin{frame}{Blocks for theorems}
	Beamer also uses blocks by default to wrap theorems.
	\begin{theorem} This is a theorem. \end{theorem}
\end{frame}


\begin{sidepic}{assets/uic_seo.jpg}{Side-Picture Slides}
	\framesubtitle{Subtitle here}
	This Beamer theme also provides a \texttt{sidepic} environment which is a wrapper on top of the \texttt{frame} environment.
	\begin{itemize}
		\item It has an optional \texttt{image} argument which will help you achieve a layout of this type.
	\end{itemize}
\end{sidepic}

\renewcommand{\algorithmicrequire}{\textbf{Input:}}
\renewcommand{\algorithmicensure}{\textbf{Output:}}

\begin{frame}{Pseudocode Example}
	\begin{algorithm}[H]
		\fontsize{8}{1}\selectfont
		\caption{Bellman-Kalaba (adapted from \hrefcol{https://www.ctan.org/tex-archive/macros/latex/contrib/algorithmicx/algorithmicx.pdf}{algorithmicx documentation}).}
		\begin{algorithmic}
			\Require $G$, $u$, $l$, $p$
			% \ENSURE optional output
			\ForAll {$v \in V(G)$}
			\State $l(v) \leftarrow \infty$
			\EndFor
			\State $l(u) \leftarrow 0$
			\Repeat
			\For {$i \leftarrow 1, n$}
			\State $min \leftarrow l(v_i)$
			\For {$j \leftarrow 1, n$}
			\If {$min > \Call{Edge}{v_i, v_j} + l(v_j)$}
			\State $min \leftarrow \Call{Edge}{v_i, v_j} + l(v_j)$ \Comment{Example comment}
			\State $p(i) \leftarrow v_j$
			\EndIf
			\EndFor
			\State $l’(i) \leftarrow min$
			\EndFor
			\State $changed \leftarrow l \not= l’$
			\State $l \leftarrow l’$
			\Until{$\neg changed$}
			% \RETURN optional return
		\end{algorithmic}
	\end{algorithm}
\end{frame}

\begin{frame}[fragile]{Typesetting Algorithms}
	\framesubtitle{Common Issues}
	\begin{itemize}
		\item There are \hrefcol{https://tex.stackexchange.com/questions/229355}{multiple environments} for typesetting algorithms.
		      \begin{itemize}
			      \item This template uses \texttt{algorithmicx} via \texttt{algpseudocode}.
			            \begin{itemize}
				            \item \texttt{algpseudocode} is (one of the) layouts for \texttt{algorithmicx}.
				            \item Others exist, as mentioned in the \hrefcol{https://www.ctan.org/tex-archive/macros/latex/contrib/algorithmicx/algorithmicx.pdf}{documentation, section 2.1}.
			            \end{itemize}
			      \item \texttt{algorithmic} is also widely used, which is an alternative to using \texttt{algorithmicx} (via a layout package).
		      \end{itemize}
		\item While \texttt{algpseudocode} tries to be close to \texttt{algorithmic}, they use different capitalization schemes
		      \begin{itemize}
			      \item \texttt{algorithmic} commands are fully capitalized (e.g. \verb|\STATE|, \verb|\FORALL|)
			      \item \texttt{algpseudocode} commands are not (e.g. \verb|\State|, \verb|\ForAll|).
			      \item \textbf{If moving an algorithm from one to the other, you will need to fix the capitalization, or the compiler will run into problems.}
		      \end{itemize}
	\end{itemize}
\end{frame}

\begin{frame}[fragile]{Typesetting Algorithms}
	\framesubtitle{More control}
	\begin{itemize}
		\item If you must use \texttt{algorithmic} instead of \texttt{algpseudocode}, some commands will need to be manually defined
		      \begin{itemize}
			      \item For example the function \verb|\Call| (or rather, the \verb|\CALL|) command.
			            \begin{itemize}
				            \item Do it as \verb|\newcommand*\CALL[2]{\textsc{#1}(#2)}|
				            \item Remember that not all fonts have a small caps variant needed by \verb|\textsc|
				            \item You can look at \texttt{uicfont.sty} to see how we augment our font with another small caps font (i.e. \textbf{PlayFair Display SC})
			            \end{itemize}
			      \item Additionally, the \verb|\Comment| of \texttt{algpseudocode} might look better than the \verb|\COMMENT| of \texttt{algorithmic}
			            \begin{itemize}
				            \item You can define a custom one to mimic \verb|\Comment|
				            \item Do it as \verb|\newcommand*\ANNOTATE[1]{\hfill\(\triangleright\) #1}|
			            \end{itemize}
		      \end{itemize}
	\end{itemize}
\end{frame}

\begin{frame}{Good luck with your presentation!}
	\begin{itemize}
		\item This template is hosted on GitHub
		      \begin{itemize}
			      \item  \texttt{https://github.com/SaverioMonaco/DESY-SINTEF}
		      \end{itemize}
		\item  I would appreciate contributions of all sorts (pull requests, identifying issues, etc).
		\item  If you have any suggestions,
		      \hrefcol{mailto:saverio.monaco@desy.de}{send them to me!}
	\end{itemize}
\end{frame}

\themecolor{dark}
\begin{frame}
	\begin{center}
		{\Huge\bfseries Thank you!}
	\end{center}
\end{frame}
\themecolor{light}

\end{document}
